% !PRINTED none
% Section-opening page: no running page number visible in the scan; part title, section
% heading, and body text follow.

\begin{center}
\gesperrt{Twelfth Section.}

\vspace{1em}
\gesperrt{Theory of the Roots of Unity.}

\vspace{0.5em}
\noindent\rule{4cm}{0.4pt}

\vspace{1em}
\S.\,131.

\vspace{0.5em}
\gesperrt{The Roots of Unity.}
\end{center}

\vspace{1em}
By a root of unity one understands in general a real or
imaginary number of such a nature that some power of it with
an integral exponent is equal to 1. We have already become
acquainted with these roots of unity in \S.\,33 and represented
them by trigonometric and exponential functions. We have,
moreover, made use of special cases, e.\,g. in \S.\,35 the third
roots of unity, which were there expressed without trigonometric
functions.

In all deeper investigations of algebraic equations a more
precise knowledge of the roots of unity and of their properties
is now indispensable. In this section we shall therefore occupy
ourselves more thoroughly with the elementary part of the
algebraic theory of the roots of unity, without making use of
their representation by trigonometric functions. On account of
the geometric application to the construction of the regular
polygons, to which we have already pointed in \S.\,33, the theory
of the roots of unity is also called \gesperrt{the theory of the division
of the circle}. It is in essence a creation of \gesperrt{Gauss} $^1$).

\noindent\rule{4cm}{0.4pt}

$^1$) \gesperrt{Gauss}, Disq. ar. Sectio VII.
